% sec6-bands.tex -- Section 6 of the C3 paper (draft 2026-11-08): result F and the F_3 target.
% Sources (run B 20261104-0759): claude-20261104-c2-direct-nd/01-DIRECT-ND.md (c2: Lemmas 1-5, (A), (I), data) and
% 02/04/06/07 (consequences by band); codex-20261104-x1-syzygies/03 (L'), 08 (S(4) independent check), 09, 10, 20-22
% (bands); x2 12-14; notes-20261104/referee-runB-direct-route.md (cold referee); Astra
% codex-20261104-runB-adjudication/00 (C2), 01-C2-GLOBAL-AUDIT.txt, 05 (S(4) numbers). Lean: K100 (run 1627).

\section{Fixed-degree bands and the \texorpdfstring{$\mathbb F_3$}{F3} target}\label{sec:bands}

This section gives a third mechanism, which works degree by degree.  It reduces
$\mathrm{ND}_d$ to a statement about the homogeneous ideals generated by row sums and by
column sums, and proves that statement for every $d<p^2$.  It cannot give growth in $\ell$ at
fixed $p$, and it is not a third proof of it; what it gives is exactness at
$v_p(k)\le2$ from a slightly lower threshold, and the $\mathbb F_3$ target, on inputs that
differ from those of Sections~\ref{sec:growth} and~\ref{sec:periodicity}.

\begin{theorem}[F]\label{thm:bands}
Let $p$ be prime, $k\ge1$, $q=p^{v_p(k)}$, and $1\le D<\min(q,p^2)$.  Over every field of
characteristic $p$, $\mathrm{onto\text{-}FPHP}^{n+k}_n$ has no PC refutation of degree at
most $D$ for any $n\ge2D-1$.  Consequently, if $v_p(k)=2$, its PC degree is exactly $p^2$
for every $n\ge2p^2-3$.
\end{theorem}

\begin{corollary}[the $\mathbb F_3$ target]\label{cor:f3}
If $9\mid k$ and $n\ge7$, $\mathrm{onto\text{-}FPHP}^{n+k}_n$ has no PC refutation of degree
at most $4$ over $\mathbb F_3$.
\end{corollary}

The corollary is the case $p=3$, $D=4$.  It has two proofs: through the identity (I) below in
degrees $d\le4$, which for $p=3$ needs only the first two bands and the connectivity theorem;
and through exact computations on three square boards (Section~\ref{sec:bands-f3}).

\paragraph{Tier and inputs.}  T-WRITTEN.  The ideal-identity proof uses the connectivity
theorem of Bj\"orner, Lov\'asz, Vre\'cica and \v Zivaljevi\'c for chessboard complexes, as stated by
Shareshian and Wachs~\cite[Thm.~1.1]{shareshianwachs}.  Degrees $d<2p$ need nothing else.
Degrees $2p\le d<p^2$ (odd $p$) also use Yal\c c\i n's theorem on equivariant Moore spaces and
the Dade group~\cite[Def.~1.1, Props.~3.8, 5.7]{yalcin2016}, Proposition~3.3 of
Miller~\cite{miller2024}, and uniqueness of Dade caps.  The input we consider most exposed is
Miller's Proposition~3.3 at its interface with the truncated chessboard cochain complex,
where the output needed is an equality of integer dimensions of Brauer quotients
(Section~\ref{sec:bands-high}).  The second proof of Corollary~\ref{cor:f3} replaces
connectivity by Church--Ellenberg's Corollary~B and exact computation.

\subsection{From ND to a kernel statement}\label{sec:bands-reduction}

We use the notation of Section~\ref{sec:defs-written} on a board with $m$ pigeons and $n$
holes.  Let $V_d$ be the span of matchings of size exactly $d$.  The \emph{labels} at level
$d$ are the pairs $(N,\mathrm{row}\ i)$ and $(N,\mathrm{col}\ j)$ with $|N|=d-1$ and $i$, $j$
free in $N$.  Let $a_d$ send a label to the top part of its relation
($\sum_jx_{N\cup\{(i,j)\}}$, resp.\ $\sum_ix_{N\cup\{(i,j)\}}$), a map from the label space
to $V_d$, and let $\beta$ send a label to its base $x_N$.

\begin{lemma}\label{lem:nd-ker}
$W_d\cap F_{d-1}=W_{d-1}+\beta(\ker a_d)$.  Hence $\mathrm{ND}_d$ holds if and only if
$\beta(\ker a_d)\subseteq W_{d-1}$.
\end{lemma}

\begin{proof}
An element of $W_d$ is $w+\sum_gc_g\,\mathrm{rel}(g)$ with $w\in W_{d-1}$ and $g$ over the
level-$d$ labels.  Its degree-$d$ part is $a_d(c)$, so it lies in $F_{d-1}$ iff
$c\in\ker a_d$, and then it equals $w-\beta(c)$.
\end{proof}

Fix total orders on pigeons and holes.  The following elements of the label space lie in
$\ker a_d$:
\begin{itemize}
\item $rr(N;a,b)=\sum_j(N\cup\{(b,j)\},\mathrm{row}\ a)-\sum_j(N\cup\{(a,j)\},\mathrm{row}\ b)$
  for $|N|=d-2$ and free pigeons $a<b$; $cc(N;a,b)$ likewise for holes;
\item $rc(N;i,j)=\sum_{i'\ne i}(N\cup\{(i',j)\},\mathrm{row}\ i)-\sum_{j'\ne j}(N\cup\{(i,j')\},\mathrm{col}\ j)$;
\item the \emph{anchored} elements $Z^{(t)}_N$, $1\le t\le d$, $|N|=d-t$:
  \[
  Z^{(t)}_N=\sum_E\Bigl(\sum_{\substack{i\ \text{free}\\ i>\max\mathrm{pig}(E)}}(N\cup E,\mathrm{row}\ i)
     -\sum_{\substack{j\ \text{free}\\ j>\max\mathrm{hol}(E)}}(N\cup E,\mathrm{col}\ j)\Bigr),
  \]
  the outer sum over $(t-1)$-matchings $E$ on the pigeons and holes free in $N$
  ($\max\emptyset=-\infty$).  Both halves map under $a_d$ to the sum of all $d$-matchings
  containing $N$, each once.
\end{itemize}

\begin{lemma}\label{lem:lower-images}
$\beta(rr)$, $\beta(cc)$, $\beta(rc)$ lie in $W_{d-1}$, and with $s=|N|$,
\[
  \beta(Z^{(t)}_N)\equiv\Bigl(\tbinom{m-s}{t}-\tbinom{n-s}{t}\Bigr)x_N\pmod{W_{d-1}}.
\]
\end{lemma}

\begin{proof}
$\beta(rr(N;a,b))$ is the difference of the row relations at $(N,b)$ and $(N,a)$, and
similarly for $cc$; $\beta(rc(N;i,j))$ is the column relation at $(N,j)$ minus the row
relation at $(N,i)$.  For $Z$, group the $E$ by their pigeon set $A$, $|A|=t-1$; the row
coefficient is the number of free pigeons above $\max A$.  For fixed $A$,
$\sum_{\mathrm{pig}(E)=A}x_{N\cup E}\equiv x_N$ modulo $W_{d-1}$, by applying one row
relation at a time at bases of size $\le d-2$.  Since $t$-subsets of the free pigeons
correspond to pairs ($A$, a new largest element), the row half contributes
$\binom{m-s}{t}x_N$; the column half contributes $\binom{n-s}{t}x_N$.
\end{proof}

If $m-n=k$ and $t<q=p^{v_p(k)}$, then $\binom{m-s}t\equiv\binom{n-s}t\pmod p$, by
$(1+z)^k=(1+z^q)^{k/q}\equiv1+O(z^q)$ over $\mathbb F_p$.  So:

\begin{proposition}\label{prop:A-to-nd}
Let $D<q$ and suppose that for every $d\le D$
\begin{center}
(A)$(m,n,d)$: \ $\ker a_d$ is spanned by the elements $rr$, $cc$, $rc$ and $Z^{(t)}_N$.
\end{center}
Then $\mathrm{ND}_1,\dots,\mathrm{ND}_D$ hold on the board, and since $W_0=0$ also
$\mathrm{U}_D$; by Proposition~\ref{prop:nd-pc} there is no PC refutation of degree at most
$D$.
\end{proposition}

(A) is homogeneous and does not see $k$; the surplus enters only through the binomial
differences.  (A) fails off the stable range, for example at $(p,m,n,d)=(3,7,4,3)$, where
$\ker a_3$ has dimension $944$ and the listed elements span $916$ (computed twice,
independently).

\subsection{The ideal identity}\label{sec:bands-identity}

Let $\mathcal A$ be the homogeneous matching ring: basis $x_M$ over matchings,
$x_Mx_{M'}=x_{M\cup M'}$ if $M,M'$ are edge-disjoint and $M\cup M'$ is a matching, and $0$
otherwise.  (So $x_{ij}^2=0$; this is the top-degree ring, not the Boolean algebra
$\mathcal B$.)  Let $R_i=\sum_jx_{ij}$, $C_j=\sum_ix_{ij}$, $I=(R_i)$, $J=(C_j)$, and
$\sigma_t$ the sum of all $t$-matchings.  The identity in degree $d$ is
\[
  \text{(I)}(m,n,d):\qquad (I\cap J)_d=(IJ)_d+\sigma_1\mathcal A_{d-1}+\sigma_p\mathcal A_{d-p},
\]
with the last term absent for $d<p$.

\begin{lemma}\label{lem:I-to-A}
If $\min(m,n)\ge2d-1$, then (I)$(m,n,d)$ implies (A)$(m,n,d)$.
\end{lemma}

\begin{proof}
To $s\in\ker a_d$ with row part $\alpha$ and column part $\gamma$ associate
$y(s)=\sum\alpha_iR_i=-\sum\gamma_jC_j\in(I\cap J)_d$ (the products are taken in
$\mathcal A$, with the base of each label as coefficient).  Then $y(rr)=y(cc)=0$,
$y(rc(N;i,j))=x_NR_iC_j$ (because $R_ix_{ij}=0$), and $y(Z^{(t)}_N)=x_N\sigma_t$.  So every
element on the right-hand side of (I) is $y$ of a combination of $rc$'s and $Z$'s: $(IJ)_d$ is
spanned by the $x_NR_iC_j$, which vanish unless $i$ and $j$ are free in $N$, and
$\sigma_1\mathcal A_{d-1}$, $\sigma_p\mathcal A_{d-p}$ by the $x_N\sigma_1$, $x_N\sigma_p$.
(The run's first form of (I) listed all $x_N\sigma_t$; the form above, with $\sigma_1$ and
$\sigma_p$ only, is the one proved.)  By (I), subtracting suitable $rc$'s and $Z$'s leaves $s'$ with $y(s')=0$, that is
$\sum\alpha'_iR_i=0$ and $\sum\gamma'_jC_j=0$ separately.  For a fixed set of $d$ pigeons
the row-only relation is a cocycle of the augmented cochain complex of the chessboard
complex $\Delta(d,n)$ in dimension $d-2$, and the $rr$'s are its coboundaries.  For
$n\ge2d-1$ the connectivity theorem~\cite[Thm.~1.1]{shareshianwachs} gives
$\tilde H_{d-2}=\tilde H_{d-3}=0$ over $\mathbb Z$, hence $\tilde H^{d-2}=0$ over every field.
So $\alpha'$ is a sum of $rr$'s, and likewise $\gamma'$ of $cc$'s (using $m\ge2d-1$).
\end{proof}

\paragraph{Degrees below $p$.}  Let $G=S_d$ act on ordered $d$-tuples by positions.  The
coefficient spaces of $(I\cap J)_d$ identify with $V^G$, where $V=P_m\otimes P_n$ and $P_m$ is
the span of functions on injective $d$-tuples in $[m]$ that are independent of at least one
position; this identification needs no division.  On the spanning tensors of deletion
indicators, the norm $N_G=\sum_{g\in G}g$ gives $x_N\sigma_1$ when the row and column
indicators omit the same position and $x_NR_iC_j$ when they omit different positions.  Hence
\[
  N_GV=(IJ)_d+\sigma_1\mathcal A_{d-1},\qquad
  (I\cap J)_d\big/\bigl((IJ)_d+\sigma_1\mathcal A_{d-1}\bigr)\cong\hat H^0(G,V).
\]
For $d<p$, $|G|=d!$ is invertible, the Tate group vanishes, and (I) holds.

\subsection{The bands \texorpdfstring{$p\le d<p^2$}{p <= d < p^2}}\label{sec:bands-high}

We give the argument in these bands in outline.  Each step was written out in full in the
run, reviewed there, and reconstructed in Astra's audit of the local-to-global argument; the
outline below follows that audit.  It is the part of the paper whose written form is furthest
from what is printed here.
For $p\le d<p^2$ write $d=ap+r$ with $0\le r<p$; then $a<p$, $v_p(d!)=a$, and
$Q=(C_p)^a$ acting on $a$ disjoint blocks of $p$ positions is a Sylow $p$-subgroup of $S_d$.
Since $[S_d:Q]$ is prime to $p$, restriction to $Q$ followed by transfer is multiplication
by a unit on $\hat H^0$, so it suffices to show that $\hat H^0(Q,V)$ is generated by
transfers of classes that become $(p-1)!\,x_N\sigma_p$ in $\mathcal A$.
\begin{enumerate}
\item \emph{One cycle} ($a=1$, any number $r$ of fixed positions).  Projecting the augmented
  chessboard cochains of $\Delta(d,m)$ to the fixed-position cochains of $\Delta(r,m)$ has a
  kernel of free modules; connectivity gives exactness below the top; dimension shifting along
  the cyclic period shows that the actual map from functions on fixed-position $r$-tuples is
  onto on $\hat H^0$.  This covers $p\le d<2p$ for every prime, including $p=2$, $d=2,3$.
\item \emph{Full-vertex summands.}  The fixed complexes of subgroups of $Q$ are again
  chessboard complexes, tight Moore complexes by connectivity; Yal\c c\i n's theorem identifies
  $P_m$ as capped endo-permutation with an explicit Dade class.  For the tensor product
  $V=P_m\otimes P_n$ the relative terms cancel, giving $[V]=2[\Omega_Q]$; its cap is the
  minimal second syzygy $C=\ker(R^a\to R)$, $R=\mathbb F_p[x_1,\dots,x_a]/(x_i^p)$.  Dade cap
  uniqueness makes every full-vertex indecomposable summand of $V$ a copy of $C$, and $\hat H^0(Q,C)$ is
  generated by transfers from the cyclic factors.
\item \emph{Proper-vertex summands.}  Miller's Proposition~3.3, applied to the truncated
  chessboard cochain complexes, gives for every non-trivial $L\le Q$ the integer identity
  $\dim\mathrm{End}(V)(L)=(h_s(m)h_s(n))^2$ for the Brauer quotients, where $h_s$ is the top
  cohomology dimension of the $L$-fixed chessboard complex.  With Higman's criterion this forces
  every proper vertex into a coordinate hyperplane $H_i$, so those summands contribute only
  transfers from the $H_i$.
\item \emph{Induction on $a$.}  Restricted to $H_i$, $V$ is the same kind of module with $a-1$
  moving blocks and $r+p$ fixed positions; the one-cycle step allows any number of fixed
  positions, so induction on $a$ reduces everything to single $p$-cycles, where the classes are
  fixed-label indicators whose transfer to $S_d$ is $(p-1)!\,x_N\sigma_p$.
\end{enumerate}
At $d\ge p^2$ the subgroup $Q$ is no longer Sylow and the argument stops; nothing here
concerns that range.

\begin{proof}[Proof of Theorem~\ref{thm:bands}]
For $d<\min(q,p^2)$ and $m,n\ge2d-1$, (I) holds (above), hence (A)
(Lemma~\ref{lem:I-to-A}), hence $\mathrm{ND}_d$ (Lemmas~\ref{lem:nd-ker},
\ref{lem:lower-images}); Proposition~\ref{prop:A-to-nd} excludes refutations of degree
$\le D$ once $n\ge2D-1$.  All spaces are defined over $\mathbb F_p$, so the statement holds
over every field of characteristic $p$.  If $v_p(k)=2$, take $D=p^2-1$; the upper bound is
Proposition~\ref{prop:riis}.
\end{proof}

\subsection{The \texorpdfstring{$\mathbb F_3$}{F3} target by finite certificates}\label{sec:bands-f3}

The second proof of Corollary~\ref{cor:f3} avoids the ideal identity.  Two lemmas move (A)
from squares to rectangles.

\begin{lemma}[hole avoidance; from Corollary~B of \cite{churchellenberg}]\label{lem:hole-avoid}
If $n\ge2e+1$, every array in $H_e$ is a sum of arrays in $H_e$ each of which avoids some
hole.  Dually, if the restriction of $v\in V_e$ to every board with one hole removed lies in
the image of $a_e$ there, then $v$ lies in the image of $a_e$.
\end{lemma}

This is the case $|S|>2e$ of Lemma~\ref{lem:small-support}.  The threshold is sharp: exact
computation shows that the statement fails at $n=2e$ on several small boards.

\begin{lemma}[pigeon induction]\label{lem:pigeon-ind}
Order the pigeons so that the removed pigeon $\omega$ is the largest.  Then
(A)$(m-1,n,d)$, (A)$(m-1,n-1,d-1)$ and Lemma~\ref{lem:hole-avoid} on $(m-1,n)$ at level $d-1$
imply (A)$(m,n,d)$.  The mirror statement removes a hole.
\end{lemma}

\begin{proof}
Split the labels of the board $(m,n)$ into those avoiding $\omega$ ($G_1$), the row labels of
$\omega$ ($G_2$), and those whose base uses $\omega$ ($G_3$), and decompose
$V_d(m)=V_d(m-1)\oplus\bigoplus_hx_{\omega h}V_{d-1}(m-1,n\setminus h)$.  Let
$s=(g_1,g_2,g_3)\in\ker a_d$.  (a) The first component of $a_d(s)$ is $a_d(g_1)$ on the board
$m-1$, so $g_1$ is in that kernel; each listed element of the smaller board is the
$G_1$-part of the element of the same name on the larger board (the extra terms involve
$\omega$ only in $G_2$, $G_3$, and $\omega$ largest preserves the order conditions of $Z$).
Subtract.  (b) If $g_1=0$, the component at $h$ reads
$g_2|_{\text{avoid }h}+a_{d-1}(g_{3,h})=0$, so by Lemma~\ref{lem:hole-avoid} $g_2$ lies in
the image of $a_{d-1}$ on $(m-1,n)$; the elements $rr(N;a,\omega)$ and $rc(N;\omega,j)$ have
$G_1$-part $0$ and $G_2$-parts exactly the needed extension sums.  Subtract.  (c) If
$g_1=g_2=0$, each $g_{3,h}$ lies in $\ker a_{d-1}$ on $(m-1,n\setminus h)$, and
$x_{\omega h}$ times a listed element at base $N'$ is the listed element of the same name at
base $N'\cup\{(\omega,h)\}$.
\end{proof}

Starting from (A) at level $1$ (which holds on every board), Lemma~\ref{lem:pigeon-ind} and
its mirror give (A) on all boards with $m,n\ge2d-1$ from the single square statement
$S(d)=$ (A)$(2d-1,2d-1,d)$.  For $p=3$ and $D=4$ the squares are $3\times3$, $5\times5$ and
$7\times7$.  $S(2)$, $S(3)$ and $S(4)$ were verified over $\mathbb F_2$ and $\mathbb F_3$ by
exact rank computations, by three implementations written by different agents; at $S(4)$ the
top matrix has rank $27{,}068$ on $58{,}800$ labels, and the $47{,}972$ listed elements span its
$31{,}732$-dimensional kernel.  With Proposition~\ref{prop:A-to-nd} this proves
Corollary~\ref{cor:f3} at tier T-WRITTEN with T-MEASURED certificates.

\paragraph{Lean.}  The library contains a conditional Lean theorem for the same target
(published as finding K100, DOI 10.5281/zenodo.24056819): if $9\mid k$, $n\ge7$, and the
spanning statement (A)$(n+k,n,r)$ holds for $r=2,3,4$, then there is no standard PC
refutation over $\mathbb F_3$ of degree at most $4$; (A) at level $1$ is proved on every board with positive numbers of pigeons and holes.\lean{\texttt{Programme.\guillemotleft 20261104-1627-proof-complexity\guillemotright}:
RemainingSpanning.no\_stdpc\_le; level 1: MatchingLabels.spanning\_one; the
semantic form ConcreteDegreeFour.no\_ppc}
The Lean theorem does not formalise Lemma~\ref{lem:pigeon-ind} or
Lemma~\ref{lem:hole-avoid}, so it needs (A) on every board $(n+k)\times n$ rather than on
three squares.  On the smallest board with positive surplus divisible by $9$, $16\times7$, (A) is verified by exact
computation at levels $2$ (kernel dimension $343$) and $3$ (kernel dimension $24{,}780$, two
independent programs); level $4$ was not decided within the time limit.

\paragraph{Who did what.}  All in run 20261104-0759 unless stated; all reviews are model
reviews.  The reduction (Lemmas~\ref{lem:nd-ker}--\ref{lem:I-to-A},
Proposition~\ref{prop:A-to-nd}, Lemma~\ref{lem:pigeon-ind}) and the form of (I) are by a Claude
agent (c2), with a Codex agent (x1), and were reviewed by a second Claude agent (c1).  The
hole-avoidance lemma is x1's, checked by c2 and adjudicated by a second Codex agent (x2).  The
identity (I) was proved band by band: in degree $2$ and whenever $d!$ is invertible by x1
(reviewed by c2); at $d=p$ by x2 (reviewed by x1 and c2); for $p<d<2p$ by x1 (reviewed by c2,
accepted by x2); in the higher bands by x1 and x2 together, with c2 reviewing each step.  The
finite certificates $S(2)$--$S(4)$ were computed by c1 and, independently, by x1, and
regenerated after the run by Astra with a third implementation.  A model referee with no
contact with the run re-derived Lemmas~1--5 and (I) for $d<2p$, re-read the sources of the
higher bands, and checked (I) and the spanning statement on ten boards with its own code.
Astra's global audit reconstructed the local-to-global argument of
Section~\ref{sec:bands-high}.  The Lean theorem is from run 20261104-1627 and was re-proved by a
second agent there.
